\section{Related Work}
\label{sec:related_work}
\subsection{Heterogeneous Workflow Scheduling}
HEFT combines task prioritization with resource selection based on estimated
finish times, while PEFT uses optimistic downstream costs to improve
placement decisions \cite{topcuoglu2002heft,arabnejad2014peft}. These methods
remain relevant as low-overhead references and initial-solution generators.
Our implementation fixes the evaluator's per-node order to a canonical
topological replay and does not optimize task order or insert tasks into idle
gaps. Accordingly, the HEFT/PEFT-style baselines in our experiments should be
understood as adaptations to this deployment evaluator, rather than a claim
to reproduce every scheduling freedom of their original formulations.

\subsection{Reliability and Hybrid Fault Tolerance}
RA-MOMA studies reliability-aware multi-objective workflow scheduling in a
multi-cloud setting using memetic search \cite{qin2023ramoma}. RECR-D addresses
fault-tolerant cost-efficient scheduling for energy- and deadline-constrained
IoT workflows in an edge--cloud continuum \cite{taghinezhad2025recrd}. HFWPF
considers hybrid fault tolerance with performance-fluctuating cloud resources
\cite{ren2026hfwpf}. These studies motivate the use of selective replication,
reexecution, and problem-specific placement decisions. They also differ in
energy modeling, resource elasticity, execution-time uncertainty, reliability
objectives, and online behavior.

The present study focuses on a different constraint: recovery must fit both
an explicit delay limit and the deadline slack left by the chosen nominal
deployment. We adapt the placement and protection mechanisms of selected
methods to a common static evaluator and describe the removed or remapped
features in Section~\ref{sec:baselines}. In particular, an adapted method's
behavior under our recovery-delay gate cannot establish that the original
method is inferior on its original problem. This distinction is essential
when the adaptation removes an uncertainty model or an online decision stage
that is central to the source method.

\subsection{Large Neighborhood Search and Learned Guidance}
ALNS extends large neighborhood search by adjusting the selection weights of
multiple destroy and repair operators according to their observed search
outcomes \cite{ropke2006alns,pisinger2007general}. This makes it suitable for
combining placement, replica, and protection-mode changes without enumerating
the entire deployment space. Our solver adds an explicit feasibility-first
phase and uses full scenario verification selectively when considering
feasible incumbents.

GCNs learn node representations through normalized graph propagation
\cite{kipf2017gcn}. Here, the resource graph represents communication-based
proximity among continuum nodes. The learned score estimates the preference
of a task--node assignment under solver demonstrations; it is not an
independent proof of feasibility or an estimate of a globally optimal
assignment. Training is supervised rather than reinforcement-based, and the
score table remains fixed throughout the search for a given instance. The
research question is whether this offline prior improves an already
constraint-aware search under independent evaluation.

\subsection{Workflow Data and External Validity}
WfCommons provides workflow instances and generators that support research
with realistic application characteristics \cite{coleman2022wfcommons}.
The experiments specified in this manuscript currently use custom synthetic
generators whose schemas and broad workflow shapes are inspired by
established scientific workflows and cluster-resource traces. They do not use
real WfCommons executions or raw Alibaba machine traces. Existing parsers and
planned data integration are therefore implementation capabilities, not
evidence that the present experimental instances are trace-driven.
